\documentclass[8pt]{extarticle} 

% \usepackage[fleqn]{amsmath}

\usepackage[margin=1in]{geometry}
\usepackage{amsmath}
% \usepackage{amsmath,amsfonts,amsthm,bm}
\usepackage{amsthm}
% \usepackage{breqn}
\usepackage{mathtools}
% \usepackage{amssymb}
% \usepackage{bbm}
\usepackage{tikz}
\usepackage[ruled,vlined,linesnumbered,lined,boxed,commentsnumbered]{algorithm2e}
\usepackage{algpseudocode}
\usepackage{siunitx}
% \usepackage{graphicx}
\usepackage{subcaption}
%% \usepackage{datetime}
\usepackage{multirow}
\usepackage{multicol}
% \usepackage{mathrsfs}
\usepackage{fancyhdr}
\usepackage{fancyvrb}
\DefineVerbatimEnvironment{Code}{BVerbatim}{baseline=t}

\usepackage{parskip} %turns off paragraph indent
\pagestyle{fancy}

\usepackage{xcolor}
\usepackage{mdframed}

\usepackage[small]{titlesec}

\usepackage{hanging}

\usepackage{listings}% http://ctan.org/pkg/listings
\lstset{
  basicstyle=\ttfamily\small,
  mathescape
}

\usepackage{tabularray}

\usetikzlibrary{positioning}
\usetikzlibrary{shapes,arrows}
% Define block styles
\tikzset{
    block/.style={rectangle, draw, line width=0.5mm, black, text width=5em, text centered,
                  rounded corners, minimum height=2em},
    line/.style={draw, -latex}
}% <- if you insist in using this in the document add this % here.

\DeclareMathOperator*{\argmin}{argmin}
\newcommand*{\argminl}{\argmin\limits}

% \newcommand{\mathleft}{\@fleqntrue\@mathmargin0pt}
% \newcommand{\R}{\mathbb{R}}
% \newcommand{\Z}{\mathbb{Z}} 
% \newcommand{\N}{\mathbb{N}}
% \newcommand{\ppartial}[2]{\frac{\partial #1}{\partial #2}}
% \newcommand{\p}{\partial}
% \newcommand{\te}[1]{\text{#1 }}
% \newcommand{\norm}[1]{\|#1\|}

\newcommand*{\vrectangle}{{\setlength{\fboxsep}{0pt}\fbox{\phantom{l}}}}

\setcounter{MaxMatrixCols}{20}

% remove excess vertical space for align equations
\setlength{\abovedisplayskip}{0pt}
\setlength{\belowdisplayskip}{0pt}
\setlength{\abovedisplayshortskip}{0pt}
\setlength{\belowdisplayshortskip}{0pt}

\newtheorem{theorem}{Theorem}[section]
\newtheorem{corollary}{Corollary}[theorem]
\newtheorem{lemma}[theorem]{Lemma}

\begin {document}

\lhead{Notes - Compiler Construction, Bill/Yuan Liu}
  
\begin{multicols*}{2}

  Gathered notes from various references:
  \begin{itemize}
  \item Modern Compiler Implementation in C \cite{appelbook}
  \item High Performance Compilers for Parallel Computing \cite{wolfebook}
  \item The implementation of Functional Programming Languages \cite{peyton1987}
  \item Modern Compiler Design \cite{grune2012}
  \item MLIR \cite{mlir2021}, \cite{lei2022}
  \end{itemize}

  \section{Data Flow Analysis}

  \subsection{Available Expression}

  purpose: compute set of expressions evaluated and not overwritten on all paths from the entry to the target program point
  
  domain: set of non-trivial expressions

  direction: forward

  join operator: $\wedge$ (set intersection), eg: must-analysis

  lattice values: $\bot$ = $\emptyset$ (empty set), $\top$ = universal set

  partial order: set_1 $\leq$ set_2 $\iff$ set_1 $\subseteq$ set_2

  transfer function for a basic block B:

  $out(B) = (in(B) - kill(B)) \cup gen(B)$

  note: kill before gen

  kill(B) = all expressions where any operands in the expression are referenced in assignment (LHS) in B

  gen(B) = assignment of an expression to a LHS variable where all operands of the expression are not redefined subsequently in B (which would make expression stale)

  boundary condition:
  
  in(entry) = $\bot$ since no expressions are computed/available at the start

  out(B) = $\top$ for $\forall B \neq entry$: join operator narrows the candidates in algorithm

  extra notes:

  use of this analysis in common subexpression elimination IR transformation

  \vfill\null
  \columnbreak
    
  \section{Liveness Analysis}

  purpose: compute variable at a program point to see wether if it is used by a code path later on

  domain: set of variables

  direction: backward

  join operator: $\cup$ (set union), eg: may-analysis

  eg: $out(s) = \cup_y in(y)$ where $y \in$ successor($s$)

  lattice values: $\bot = \emptyset$ (empty set of variables), $\{var_i, ..\}$, $\top$ = universal set of variables

  partial order: $\leq = \subseteq$

  transfer function:

  $out(B) = (in(B) - kill(B)) \cup gen(B)$

  note: apply kill set first before applying gen set

  $gen(s) = use(s) = \{ x:\ if\ S: lhs = .. x .. \}$ = all variables referened in RHS of an assignment

  $kill(s) = def(s) = \{ x:\ if\ S: x = ... \}$ = LHS variable in an assignment

  boundary condition:

  $out(exit) = \emptyset$ if pure local analysis, or

  $out(exit)$ = set of variables referenced by return statement
  
  \vfill\null
  \columnbreak

  \section{Reaching Definition May-Analysis}

  purpose: compute if there is at least 1 path incoming to the current program point has the definition $D: x = ...$ and there exists no reassignment of $x$ with something else after $D$
  
  domain: set of definition sites
  
  direction: forward

  join operator = $\cup$ (set union)

  transfer function:

  $out(B) = (in(B) - kill(B)) \cup gen(B)$
  
  kill(B) = existing definitions overwritten (involving variable being defined/assigned) in B

  gen(B) = locally exposed new definitions in basic block B

  bounary condition:

  $out(entry) = \emptyset$

  \section{Reaching Definition Must-Analysis}

  purpose: compute if all paths incoming to the current program point has the definition $D: x = ...$ and there exists no reassignment of $x$ with something else after $D$
  
  join operator = $\wedge$ (set intersection)

  everything else is the same as reaching definition may-analysis

  \vfill\null
  \columnbreak

  \section{Constant Propagation, IR Rewrite}
  
  purpose: computes where a variable is known to be a compile time known value c, or undefined, or dynamic/undeterminable value

  domain: mapping of variable to abstract value

  direction: forward

  lattice values: $\bot$, $c_i \in Constant$, $\top$

  partial ordering: $\bot < c_i < \top$
  
  join operator:

  \begin{verbatim}
  f(a,b) = case (a,b) of
             (top, _) => top // absorb; has priority over (bot,v)
             (bot, v) => v // bot acts as identity for join
             (c, c) => c // same constant in different branches
             (c1, c2) | c1 != c2 => top // dynamic/over-determined
  \end{verbatim}

  boundary condition:

  \begin{verbatim}
  in(entry)(v) = top if v is a parameter/external unknown value
               = bot if v is uninitialized

  out(B)(v) = bot for all v, all B != entry
  \end{verbatim}

  transfer function: not applicable using generic DFA equation

  update mapping of variable to abstract value for specific statements:

  \begin{verbatim}
  out = in + mapping
  \end{verbatim}

  statement with binary operator S: $x = op(a, b)$
  \begin{verbatim}
  f(a, b) = case (a,b) of
              (top, _) => top // override; highest precedence
              (bot, v) => bot // still waiting for an operand
              (c1, c2) => op(c1, c2) //both operands are constants

  mapping[x] = f(a, b)
  \end{verbatim}

  statement with assignment operator $S: x = y$

  \begin{verbatim}
  f(x) = mapping[y]
  \end{verbatim}
  
  external variables, eg:
  \begin{verbatim}
    x = read_io() //dynamic
  \end{verbatim}, then
  $x$ is mapped to $\top$

  other statements that do not modify variables, then use identity mapping:
  \begin{verbatim}
  out = in
  \end{verbatim}

  transfer mapping composition for basic block $B$: compose statement-level mappings within $B$:
  \begin{verbatim}
  out(B) = f_n(...(f_0(in(B))))
  \end{verbatim}  
  
  this transformation still shares worklist based approach typical in DFA framework, using pseudo-topological ordering or reverse post order (reverse of post-visit DFS order)
  
  \vfill\null
  \columnbreak

  \section{Common Subexpression Elimination, IR Rewrite}
  
  dependent analysis:

  available expression analysis output $in(B)$ computed for every $B$ in the program

  direction: forward

  lattice values: set of available expressions

  partial ordering:

  $set_1 \leq set_2 \iff set_1 \subseteq set_2$

  bookkeeping:

  expression table lookup mapping operands tuple of expression to variable (SSA) holding the result (LHS of assignment to expressions); this expression is queried and mutated during  traversal of graph by the algorithm

  traversal order: dominator tree order (eg: dfs over dominator tree)

  algorithm:
  
  \begin{verbatim}
  f(B) =  
    // in[B] computed by available expr analysis
    // seed with incoming paths' expressions
    // as the set of available expressions
    local_avai_exprs = in[B].clone()

    // for rollback at the exit of the scope
    keys_inserted_in_scope = {}

    for instr S in B in sequence:
      let (lhs = rhs) = S

      // expression under consideration
      // eg: a + b if original assignment is ..=a+b
      key = rhs

      // check using local_avai_exprs to make sure
      // expression is available on every incoming paths
      //
      // use of exprs_table to see dominating variable
      if key in local_avai_exprs and
         key in exprs_table:
        var_existing = exprs_table.lookup(key)
        replace S with `lhs = var_existing`
      else:
        exprs_table.insert(key, lhs)
        local_avai_exprs.insert(key)
        keys_inserted_in_scope.insert(key)

    // recursion using immediate domination relation
    // this guarantees items processed in exprs_table
    // is seen when we visit dominated nodes
    for child in dominator_tree[B]:
      f(child)

    for k in keys_inserted_in_scope:
      exprs_table.erase(k)
  \end{verbatim}
  
  \vfill\null
  \columnbreak

  \section{Copy Propagation, IR Rewrite}
  
  todo

  \vfill\null
  \columnbreak

  \section{Dead Code Elimination, IR Rewrite}
  
  todo

  \vfill\null
  \columnbreak

  \section{Partial Redundancy Elimination, IR Rewrite}
  
  defer for now

  \vfill\null
  \pagebreak

  \section{MFP}
  todo
  \section{MOP}
  todo
  
  \section{Dominance}
  todo
  
  \section{Immediate Dominance}
  todo
  
  \section{Dominance Frontier}
  todo
  
  \section{Phi-Function Placement}
  todo
  
  \section{SSA Renaming}
  todo
  
  \section{Region Analysis}
  todo
  
  \section{Loop Inversion}
  transformation of a pre-test loop into a post-test loop
  
\bibliographystyle{plain}
\bibliography{notebook}
    
\end{multicols*}

\end {document}
